\documentclass[10pt]{beamer}
\usetheme{metropolis}
\usepackage{graphicx}
\usepackage{listings}
\usepackage{booktabs}
\usepackage{xcolor}
\definecolor{codebg}{HTML}{F5F5F5}
\definecolor{codegreen}{HTML}{2E7D32}
\definecolor{codegray}{HTML}{757575}
\definecolor{codeblue}{HTML}{1565C0}
\lstdefinestyle{rstyle}{language=R,backgroundcolor=\color{codebg},basicstyle=\ttfamily\scriptsize,
  keywordstyle=\color{codeblue}\bfseries,stringstyle=\color{codegreen},
  commentstyle=\color{codegray}\itshape,breaklines=true,frame=single,
  rulecolor=\color{codegray!30},showstringspaces=false,tabsize=2,
  morekeywords={library,ggplot,aes,geom_errorbar,geom_pointrange,geom_ribbon,geom_smooth,
    geom_linerange,geom_crossbar,stat_summary,mean_cl_normal,mean_se,
    scale_color_manual,theme_minimal,labs,after_stat}}
\lstdefinestyle{pythonstyle}{language=Python,backgroundcolor=\color{codebg},basicstyle=\ttfamily\scriptsize,
  keywordstyle=\color{codeblue}\bfseries,stringstyle=\color{codegreen},
  commentstyle=\color{codegray}\itshape,breaklines=true,frame=single,
  rulecolor=\color{codegray!30},showstringspaces=false,tabsize=4,
  morekeywords={import,from,as,pd,plt,sns,np,fig,ax,errorbar,fill_between,sem,pointplot}}
\title{Uncertainty: Error Bars, Confidence \& Gradient}
\subtitle{Workshop 3 --- Module 7}
\author{Prof.Asoc.\ Endri Raco}
\institute{Data Visualization for Data Scientists \\ UNYT Tirana}
\date{2025/26}
\begin{document}

\begin{frame}
  \titlepage
\end{frame}

\begin{frame}{Module Roadmap}
  \begin{enumerate}
    \item Why visualize uncertainty: point estimates are misleading
    \item Three error bar types: SE, 95\% CI, SD
    \item Pointrange: the better error bar chart
    \item Confidence ribbons: uncertainty on time series
    \item Gradient density bands: continuous probability
    \item Forest plots: meta-analysis style
  \end{enumerate}
  \begin{exampleblock}{Learning Outcome}
    You will always show uncertainty alongside point estimates, choose
    the correct error metric (SE vs CI vs SD), and produce
    publication-quality uncertainty visualizations in R and Python.
  \end{exampleblock}
\end{frame}

\section{Why Show Uncertainty}

\begin{frame}{Point Estimates Alone Are Misleading}
  \begin{center}
    \includegraphics[width=0.92\textwidth]{figures_w03m07/why_uncertainty}
  \end{center}
  \begin{alertblock}{Pro-Tip}
    A bar chart of means without error bars implies precision that does
    not exist. Overlapping 95\% CIs indicate the difference is likely
    not statistically significant. \textbf{Every mean needs an uncertainty
    interval.}
  \end{alertblock}
\end{frame}

\section{Error Bar Types}

\begin{frame}{SE vs 95\% CI vs SD}
  \begin{center}
    \includegraphics[width=0.88\textwidth]{figures_w03m07/error_bar_types}
  \end{center}
  \begin{exampleblock}{Data Insight}
    \textbf{SE} (standard error) = precision of the mean estimate.
    \textbf{95\% CI} ($\pm 1.96 \cdot SE$) = range likely containing the
    true mean. \textbf{SD} = spread of individual observations. Use CI
    for inference; SD for describing variability. Always label which one
    you are showing.
  \end{exampleblock}
\end{frame}

\begin{frame}[fragile]{Error Bars in R vs Python}
  \begin{columns}[T]
    \begin{column}{0.48\textwidth}
      \textbf{R (ggplot2)}
\begin{lstlisting}[style=rstyle]
# Auto-computed mean + 95% CI
ggplot(apps, aes(x = Category,
    y = Rating)) +
  stat_summary(
    fun.data = mean_cl_normal,
    geom = "errorbar",
    width = 0.2,
    color = "#E53935",
    linewidth = 1) +
  stat_summary(
    fun = mean,
    geom = "point",
    size = 3,
    color = "#1565C0") +
  coord_flip() +
  theme_minimal()

# mean_cl_normal  -> 95% CI
# mean_se         -> +/- 1 SE
# mean_sdl(mult=1)-> +/- 1 SD
\end{lstlisting}
    \end{column}
    \begin{column}{0.48\textwidth}
      \textbf{Python}
\begin{lstlisting}[style=pythonstyle]
from scipy.stats import sem

# Compute manually
grouped = apps.groupby("Category")
stats = grouped["Rating"].agg(
    ["mean", "count", "std"])
stats["se"] = stats["std"] / \
    np.sqrt(stats["count"])
stats["ci95"] = 1.96 * stats["se"]

# Plot
ax.errorbar(
    stats["mean"],
    range(len(stats)),
    xerr=stats["ci95"],
    fmt="o",
    color="#1565C0",
    markersize=6,
    capsize=5,
    linewidth=1.5)

# seaborn (auto CI):
sns.pointplot(data=apps,
    x="Rating", y="Category",
    errorbar=("ci", 95))
\end{lstlisting}
    \end{column}
  \end{columns}
\end{frame}

\section{Pointrange}

\begin{frame}{Pointrange: The Better Bar Chart}
  \begin{center}
    \includegraphics[width=0.72\textwidth]{figures_w03m07/pointrange}
  \end{center}
  \begin{alertblock}{Pro-Tip}
    Replace bar charts of means with \textbf{pointrange} plots: a point
    for the mean + horizontal lines for the CI. This eliminates the
    misleading bar area and forces the viewer to focus on the estimate
    and its precision.
  \end{alertblock}
\end{frame}

\begin{frame}[fragile]{Pointrange in Code}
  \begin{columns}[T]
    \begin{column}{0.48\textwidth}
      \textbf{R}
\begin{lstlisting}[style=rstyle]
ggplot(apps, aes(
    x = fct_reorder(Category,
      Rating, .fun = mean),
    y = Rating)) +
  stat_summary(
    fun.data = mean_cl_normal,
    geom = "pointrange",
    color = "#1565C0",
    size = 0.5,
    linewidth = 1) +
  coord_flip() +
  geom_hline(
    yintercept = mean(
      apps$Rating, na.rm = TRUE),
    linetype = "dashed",
    color = "grey50") +
  theme_minimal() +
  labs(title = "Mean Rating (95% CI)",
       x = NULL, y = "Rating")
\end{lstlisting}
    \end{column}
    \begin{column}{0.48\textwidth}
      \textbf{Python}
\begin{lstlisting}[style=pythonstyle]
fig, ax = plt.subplots()

# Sorted by mean
stats_sorted = stats.sort_values(
    "mean")

ax.errorbar(
    stats_sorted["mean"],
    range(len(stats_sorted)),
    xerr=stats_sorted["ci95"],
    fmt="o",
    color="#1565C0",
    markersize=7,
    capsize=5,
    linewidth=1.5,
    capthick=1.5)

ax.axvline(
    x=stats_sorted["mean"].mean(),
    color="grey", lw=0.8, ls="--")

ax.set_yticks(
    range(len(stats_sorted)))
ax.set_yticklabels(
    stats_sorted.index)
\end{lstlisting}
    \end{column}
  \end{columns}
\end{frame}

\section{Confidence Ribbons}

\begin{frame}{Ribbon: Uncertainty on Time Series}
  \begin{center}
    \includegraphics[width=0.72\textwidth]{figures_w03m07/confidence_ribbon}
  \end{center}
  \begin{exampleblock}{Data Insight}
    A ribbon shades the region between the lower and upper CI bounds
    around a trend line. Use a lighter shade for the wider 95\% CI and a
    darker shade for the narrower $\pm$1 SE. ggplot2's
    \texttt{geom\_smooth(se = TRUE)} does this automatically;
    matplotlib requires \texttt{ax.fill\_between()}.
  \end{exampleblock}
\end{frame}

\begin{frame}[fragile]{Ribbon in R vs Python}
  \begin{columns}[T]
    \begin{column}{0.48\textwidth}
      \textbf{R}
\begin{lstlisting}[style=rstyle]
# Automatic via geom_smooth
ggplot(df, aes(x = year,
               y = revenue)) +
  geom_smooth(method = "lm",
    se = TRUE,          # CI band
    color = "#1565C0",
    fill = "#1565C0",
    alpha = 0.15) +
  geom_point(color = "#1565C0") +
  theme_minimal()

# Manual ribbon
ggplot(df, aes(x = year)) +
  geom_ribbon(aes(
    ymin = lower_ci,
    ymax = upper_ci),
    fill = "#1565C0",
    alpha = 0.15) +
  geom_line(aes(y = mean),
    color = "#1565C0", lw = 1.2)
\end{lstlisting}
    \end{column}
    \begin{column}{0.48\textwidth}
      \textbf{Python}
\begin{lstlisting}[style=pythonstyle]
# Manual ribbon
ax.plot(years, means, "-o",
    color="#1565C0", lw=2)

# 95% CI band
ax.fill_between(years,
    means - 1.96 * ses,
    means + 1.96 * ses,
    alpha=0.12, color="#1565C0",
    label="95% CI")

# +/- 1 SE (darker)
ax.fill_between(years,
    means - ses,
    means + ses,
    alpha=0.25, color="#1565C0",
    label="1 SE")

ax.legend(fontsize=7)
\end{lstlisting}
    \end{column}
  \end{columns}
\end{frame}

\section{Gradient Density}

\begin{frame}{Gradient Bands: Continuous Probability}
  \begin{center}
    \includegraphics[width=0.72\textwidth]{figures_w03m07/gradient}
  \end{center}
  \begin{alertblock}{Pro-Tip}
    Gradient (faded) bands convey continuous probability: darker =
    more likely, lighter = less likely. This avoids the sharp boundary
    problem of single-CI ribbons. In R: \texttt{ggdist::stat\_lineribbon()}.
    In Python: loop \texttt{fill\_between()} with decreasing alpha at
    increasing multiples of SE.
  \end{alertblock}
\end{frame}

\section{Forest Plots}

\begin{frame}{Forest Plot: Meta-Analysis Style}
  \begin{center}
    \includegraphics[width=0.72\textwidth]{figures_w03m07/forest}
  \end{center}
  \begin{exampleblock}{Data Insight}
    Forest plots show effect sizes with CIs for multiple studies, with
    a pooled (diamond) estimate at the bottom. The vertical dashed line
    at zero (or the null value) is the reference. If a CI crosses zero,
    the effect is not significant. Standard in medical and social science
    meta-analysis.
  \end{exampleblock}
\end{frame}

\section{Chart Selection}

\begin{frame}{Uncertainty Visualization Comparison}
  \begin{center}
    \includegraphics[width=0.92\textwidth]{figures_w03m07/comparison_card}
  \end{center}
\end{frame}

\begin{frame}{When to Use Which}
  \centering
  \begin{tabular}{lll}
    \toprule
    \textbf{Scenario} & \textbf{Chart} & \textbf{Error Metric} \\
    \midrule
    Compare group means & Pointrange & 95\% CI \\
    Time series trend & Ribbon / gradient & 95\% CI or SE \\
    Describe data spread & Error bar with SD & SD \\
    Meta-analysis & Forest plot & 95\% CI per study \\
    Regression fit & \texttt{geom\_smooth(se=T)} & CI on fitted line \\
    Forecast & Gradient / fan chart & Prediction intervals \\
    \bottomrule
  \end{tabular}
\end{frame}

\section{Complete Examples}

\begin{frame}[fragile]{R: Pointrange (Publication Quality)}
  \begin{columns}[T]
    \begin{column}{0.52\textwidth}
\begin{lstlisting}[style=rstyle]
apps |>
  filter(!is.na(Rating),
    Category %in% top6) |>
  ggplot(aes(
    x = fct_reorder(Category,
      Rating, .fun = mean),
    y = Rating)) +
  stat_summary(
    fun.data = mean_cl_normal,
    geom = "pointrange",
    color = "#1565C0",
    size = 0.5,
    linewidth = 1) +
  coord_flip() +
  geom_hline(
    yintercept = grand_mean,
    linetype = "dashed",
    color = "grey50") +
  theme_minimal() +
  labs(
    title = "SOCIAL Has the
      Highest Mean Rating",
    subtitle = "Mean +/- 95% CI",
    x = NULL, y = "Rating")
\end{lstlisting}
    \end{column}
    \begin{column}{0.45\textwidth}
      \includegraphics[width=\textwidth]{figures_w03m07/r_output}
    \end{column}
  \end{columns}
\end{frame}

\begin{frame}[fragile]{Python: Confidence Ribbon}
  \begin{columns}[T]
    \begin{column}{0.52\textwidth}
\begin{lstlisting}[style=pythonstyle]
fig, ax = plt.subplots(
    figsize=(7, 5))

# Trend line
ax.plot(years, revenue, "-o",
    color="#2E7D32", lw=2,
    markersize=5)

# 95% CI ribbon
ax.fill_between(years,
    revenue - 1.96 * se,
    revenue + 1.96 * se,
    alpha=0.15, color="#2E7D32")

ax.set_title(
    "Revenue Trend with 95% CI",
    fontweight="bold")
ax.set_xlabel("Year")
ax.set_ylabel("Revenue ($M)")
ax.spines["top"].set_visible(False)
ax.spines["right"].set_visible(False)
\end{lstlisting}
    \end{column}
    \begin{column}{0.45\textwidth}
      \includegraphics[width=\textwidth]{figures_w03m07/py_output}
    \end{column}
  \end{columns}
\end{frame}

\section{Synthesis}

\begin{frame}{Key Takeaways}
  \begin{enumerate}
    \item \textbf{Never show a mean without uncertainty}. Bar charts of
          means without error bars are actively misleading.
    \item \textbf{SE} = precision of the mean; \textbf{95\% CI} = range
          containing the true mean; \textbf{SD} = data spread. Always
          label which one you show.
    \item \textbf{Pointrange} replaces the bar+errorbar pattern:
          \texttt{stat\_summary(geom = "pointrange")} (R) or
          \texttt{ax.errorbar(fmt="o")} (Python).
    \item \textbf{Ribbons} show uncertainty on time series:
          \texttt{geom\_ribbon()} or \texttt{ax.fill\_between()}.
    \item \textbf{Gradient bands} avoid sharp CI boundaries by fading
          with probability distance.
    \item \textbf{Forest plots} summarize multiple studies with a pooled
          estimate --- standard in meta-analysis.
  \end{enumerate}
\end{frame}

\begin{frame}{Essential Readings}
  \begin{itemize}
    \item Cumming, G. (2014). ``The new statistics: Why and how.''
          \emph{Psychological Science}, 25(1), 7--29.
    \item Kay, M. et al. (2016). ``When (ish) is my bus?''
          \emph{CHI '16}. (Gradient uncertainty.)
    \item Wilke, C.~O. (2019). \emph{Fundamentals of Data Visualization},
          Chapter 16 (Visualizing Uncertainty).
    \item Wickham, H. (2016). \emph{ggplot2}, Chapter 5.4
          (Statistical Summaries).
  \end{itemize}
  \vspace{6pt}
  \textbf{Next Module}: Small Multiples \& Faceted Displays (W03-M08)
\end{frame}

\end{document}
